% GENERATED FILE. Edit canonical lessons, then run npm run generate:books.
% canonical-source-hash: fnv1a64:da69e3f6ff8b4abd
% canonical-lessons: HI-C282-batana, HI-C282-dur, HI-C282-kaisejatehain, HI-R282-asking-the-way

\chapter{Asking the Way}
\label{ch:hi-asking-the-way}
\hlchaptermodality{282}

\begin{chapteropening}
\textbf{By the end of this chapter:} \emph{I can ask how to get somewhere and how far it is (sṭeśan kaise jāte haiṁ?, kitnī dūr hai?), follow the answer (sīdhe jāie, bāẽ muṛiye), and ask someone to tell me (batāie).}

Complete HI-R282-asking-the-way: ask the way, ask how far, and follow the answer.
\end{chapteropening}

\section[\texorpdfstring{\hi{बताना} (batānā)}{बताना (batānā)}]{\hi{बताना} (batānā) — to tell, to let someone know}

\label{lesson:HI-C282-batana}



\begin{quote}
Before the new one: say \textquotedblleft{}he said that he would come tomorrow\textquotedblright{}. (\emph{usne kahā ki vah kal āegā}.) \textbf{\hi{कहना}} reports words. This lesson's verb hands someone a piece of information.
\end{quote}

\subsection*{You'll want to know: \hi{बताना}}
\textbf{\hi{बताना}} — \emph{batānā} — \textquotedblleft{}to tell, to let someone know\textquotedblright{}.

\begin{quote}
\textbf{\hi{मुझे} \hi{रास्ता} \hi{बताइए।}}
\end{quote}

\begin{quote}
\emph{mujhe rāstā batāie.} — \textquotedblleft{}Please tell me the way.\textquotedblright{}
\end{quote}

The one who is told takes \textbf{\hi{को}} or its short form: \textbf{\hi{मुझे}}, \textbf{\hi{उसे}}. What is told is the thing, and in the past the verb follows it:

\begin{itemize}
\raggedright
  \item \textbf{\hi{उसने} \hi{मुझे} \hi{अपना} \hi{नाम} \hi{बताया।}} (\emph{usne mujhe apnā nām batāyā.}) — She told me her name.
  \item \textbf{\hi{क्या} \hi{आप} \hi{मुझे} \hi{बता} \hi{सकते} \hi{हैं}?} (\emph{kyā āp mujhe batā sakte haiṁ?}) — Can you tell me?
\end{itemize}

\textbf{\hi{कहना}} reports what someone said: \textbf{\hi{उसने} \hi{कहा} \hi{कि} …}. \textbf{\hi{बताना}} gives someone news or directions: \textbf{\hi{उसने} \hi{मुझे} \hi{बताया} \hi{कि} …} (\emph{usne mujhe batāyā ki …}), \textquotedblleft{}she told me that …\textquotedblright{}.

\begin{cousinweb}
\textbf{\hi{बताना}} (\emph{batānā}) has no row in CDIAL, and its origin is not settled. Its ending is the \textbf{-\hi{आना}} of \textbf{\hi{दिखाना}}, \textquotedblleft{}to show\textquotedblright{}, and \textbf{\hi{समझाना}}, \textquotedblleft{}to explain\textquotedblright{}: verbs of passing something on to someone.
\end{cousinweb}

\subsection*{Your turn}
Say these aloud:

\begin{itemize}
\raggedright
  \item \emph{batānā}, then \emph{mujhe rāstā batāie}
  \item \textquotedblleft{}She told me her name\textquotedblright{} — \emph{usne mujhe apnā nām batāyā}
  \item \textquotedblleft{}Can you tell me?\textquotedblright{} — \emph{kyā āp mujhe batā sakte haiṁ?}
\end{itemize}

\begin{tcolorbox}[breakable,colback=teal!4,colframe=teal!35!black,title={Before you move on}]
What does \textbf{\hi{बताना}} mean? (\textquotedblleft{}To tell, to let someone know\textquotedblright{}.) \textquotedblleft{}Please tell me the way\textquotedblright{}? (\textbf{\hi{मुझे} \hi{रास्ता} \hi{बताइए}}.) Which verb reports what someone said? (\textbf{\hi{कहना}}.)
\end{tcolorbox}
\section[\texorpdfstring{\hi{दूर} (dūr)}{दूर (dūr)}]{\hi{दूर} (dūr) — far}

\label{lesson:HI-C282-dur}



\begin{quote}
Before the new one: say \textquotedblleft{}near the house\textquotedblright{}. (\textbf{\hi{घर} \hi{के} \hi{पास}}.) This lesson's word is its opposite.
\end{quote}

\subsection*{You'll want to know: \hi{दूर}}
\textbf{\hi{दूर}} — \emph{dūr} — \textquotedblleft{}far, far away\textquotedblright{}.

\begin{quote}
\textbf{\hi{स्टेशन} \hi{यहाँ} \hi{से} \hi{कितनी} \hi{दूर} \hi{है}?}
\end{quote}

\begin{quote}
\emph{sṭeśan yahā̃ se kitnī dūr hai?} — \textquotedblleft{}How far is the station from here?\textquotedblright{}
\end{quote}

Hindi asks \textbf{\hi{कितनी} \hi{दूर}}, with the feminine \textbf{\hi{कितनी}}, whatever the place. \textbf{\hi{दूर}} itself never changes: \textbf{\hi{गाँव} \hi{दूर} \hi{है}}, \textbf{\hi{नदी} \hi{दूर} \hi{है}}. The place you measure from takes \textbf{\hi{से}}, as in \textbf{\hi{यहाँ} \hi{से}}.

\begin{itemize}
\raggedright
  \item \textbf{\hi{बहुत} \hi{दूर} \hi{नहीं} \hi{है।}} (\emph{bahut dūr nahīṃ hai.}) — It's not very far.
  \item \textbf{\hi{मेरा} \hi{घर} \hi{स्टेशन} \hi{से} \hi{दूर} \hi{है।}} (\emph{merā ghar sṭeśan se dūr hai.}) — My house is far from the station.
\end{itemize}

\begin{cousinweb}
\textbf{\hi{दूर}} (\emph{dūr}) comes from Sanskrit \emph{dūra}, \textquotedblleft{}distant\textquotedblright{} (CDIAL 6495), and has kept its shape almost unchanged.
\end{cousinweb}

\subsection*{Your turn}
Say these aloud:

\begin{itemize}
\raggedright
  \item \emph{dūr}, then \emph{kitnī dūr}
  \item \textquotedblleft{}How far is the station from here?\textquotedblright{} — \emph{sṭeśan yahā̃ se kitnī dūr hai?}
  \item \textquotedblleft{}It's not very far\textquotedblright{} — \emph{bahut dūr nahīṃ hai}
\end{itemize}

\begin{tcolorbox}[breakable,colback=teal!4,colframe=teal!35!black,title={Before you move on}]
What does \textbf{\hi{दूर}} mean? (\textquotedblleft{}Far\textquotedblright{}.) \textquotedblleft{}How far is it?\textquotedblright{} (\textbf{\hi{कितनी} \hi{दूर} \hi{है}?}) \textquotedblleft{}From here\textquotedblright{}? (\textbf{\hi{यहाँ} \hi{से}}.)
\end{tcolorbox}
\section[\texorpdfstring{\hi{स्टेशन} \hi{कैसे} \hi{जाते} \hi{हैं}? (sṭeśan kaise jāte haiṁ?)}{स्टेशन कैसे जाते हैं? (sṭeśan kaise jāte haiṁ?)}]{\hi{स्टेशन} \hi{कैसे} \hi{जाते} \hi{हैं}? (sṭeśan kaise jāte haiṁ?) — asking the way, and following it}

\label{lesson:HI-C282-kaisejatehain}



\begin{quote}
Say \textquotedblleft{}Please tell me the way\textquotedblright{} and \textquotedblleft{}How far is it?\textquotedblright{} (\emph{mujhe rāstā batāie}, \emph{kitnī dūr hai?})
\end{quote}

\subsection*{You'll want to know: asking the way}
\begin{quote}
\textbf{\hi{माफ़} \hi{कीजिए}, \hi{स्टेशन} \hi{कैसे} \hi{जाते} \hi{हैं}?}
\end{quote}

\begin{quote}
\emph{māf kījiye, sṭeśan kaise jāte haiṁ?} — \textquotedblleft{}Excuse me, how do I get to the station?\textquotedblright{}
\end{quote}

\textbf{\hi{जाते} \hi{हैं}} with no one named is \textquotedblleft{}people go\textquotedblright{}: it asks for the way, not for anyone's plans. Or ask where it is: \textbf{\hi{स्टेशन} \hi{कहाँ} \hi{है}?} (\emph{sṭeśan kahā̃ hai?})

The answer often comes in minutes: \textbf{\hi{पाँच} \hi{मिनट} \hi{का} \hi{रास्ता} \hi{है}} (\emph{pā̃c minaṭ kā rāstā hai}), \textquotedblleft{}it's five minutes away\textquotedblright{}.

\subsection*{You'll want to know: following the way}
\begin{quote}
\textbf{\hi{सीधे} \hi{जाइए}, \hi{फिर} \hi{बाएँ} \hi{मुड़िए।}}
\end{quote}

\begin{quote}
\emph{sīdhe jāie, phir bāẽ muṛiye.} — \textquotedblleft{}Go straight on, then turn left.\textquotedblright{}
\end{quote}

\textbf{\hi{मुड़ना}} (\emph{muṛnā}) is \textquotedblleft{}to turn\textquotedblright{}, and \textbf{\hi{मुड़िए}} asks it politely, as \textbf{\hi{बोलिए}} does. \textbf{\hi{सीधे}}, \textbf{\hi{बाएँ}}, \textbf{\hi{दाएँ}} (\emph{sīdhe}, \emph{bāẽ}, \emph{dāẽ}) are the forms of \textbf{\hi{सीधा}}, \textbf{\hi{बायाँ}} and \textbf{\hi{दायाँ}} that say which way.

\begin{itemize}
\raggedright
  \item \textbf{\hi{चौराहे} \hi{से} \hi{दाएँ} \hi{मुड़िए।}} (\emph{caurāhe se dāẽ muṛiye.}) — Turn right at the crossroads.
  \item \textbf{\hi{पुल} \hi{के} \hi{बाद} \hi{दूसरी} \hi{गली} \hi{में} \hi{बाएँ} \hi{मुड़िए।}} (\emph{pul ke bād dūsrī galī meṁ bāẽ muṛiye.}) — After the bridge, take the second side street on the left. \textbf{\hi{गली}} is feminine, so \textbf{\hi{दूसरी}}.
\end{itemize}

\subsection*{Your turn}
\begin{itemize}
\raggedright
  \item \emph{Say it:} \textquotedblleft{}Excuse me, where is the post office?\textquotedblright{} — \emph{māf kījiye, ḍākghar kahā̃ hai?}
\end{itemize}

\begin{tcolorbox}[breakable,colback=teal!4,colframe=teal!35!black,title={Before you move on}]
How do you ask the way to the station? (\textbf{\hi{स्टेशन} \hi{कैसे} \hi{जाते} \hi{हैं}?}) Then say \textquotedblleft{}go straight on\textquotedblright{} and \textquotedblleft{}turn left\textquotedblright{}. (\textbf{\hi{सीधे} \hi{जाइए}}, \textbf{\hi{बाएँ} \hi{मुड़िए}}.)
\end{tcolorbox}
\section[(review)]{Review — asking the way}

\label{lesson:HI-R282-asking-the-way}



\begin{quote}
Ask \textquotedblleft{}How do I get to the station?\textquotedblright{} (\emph{sṭeśan kaise jāte haiṁ?}) Then \textquotedblleft{}How far is it?\textquotedblright{} (\emph{kitnī dūr hai?})
\end{quote}

\subsection*{Your turn}
Say these aloud:

\begin{itemize}
\raggedright
  \item Excuse me, how do I get to the market?
  \item Is it far? No, it's ten minutes away.
  \item Go straight on, then turn right at the crossroads.
  \item Can you tell me the way to the hospital?
  \item Yesterday a man told me the way.
\end{itemize}

\begin{tcolorbox}[breakable,colback=teal!4,colframe=teal!35!black,title={Before you move on}]
\textquotedblleft{}Please tell me the way\textquotedblright{}? (\textbf{\hi{मुझे} \hi{रास्ता} \hi{बताइए}}.) \textquotedblleft{}How far?\textquotedblright{} (\textbf{\hi{कितनी} \hi{दूर}?}) \textquotedblleft{}Turn right\textquotedblright{}? (\textbf{\hi{दाएँ} \hi{मुड़िए}}.) \textquotedblleft{}A man told me the way\textquotedblright{}? (\textbf{\hi{एक} \hi{आदमी} \hi{ने} \hi{मुझे} \hi{रास्ता} \hi{बताया}}: \textbf{\hi{ने}}, and \textbf{\hi{बताया}} for \textbf{\hi{रास्ता}}.)
\end{tcolorbox}
